\section*{Chapitre \ref{ch:g5:natural-hazards} --- Les aléas naturels}

\begin{solution}{exo:g5:natural-hazards:1}
(a) Un simple \omterm{def:g5:natural-hazards:hazard}{aléa naturel}~: personne n'est touché. (b) Une \omterm{def:g5:natural-hazards:disaster}{catastrophe
naturelle}~: des logements sont endommagés. (c) Un simple \omterm{def:g5:natural-hazards:hazard}{aléa naturel}.
(d) Une \omterm{def:g5:natural-hazards:disaster}{catastrophe naturelle}.
\end{solution}

\begin{solution}{exo:g5:natural-hazards:2}
Venus de l'air~: une canicule, une \omterm{def:g5:natural-hazards:storm}{tempête}, une \omterm{def:g5:natural-hazards:drought}{sécheresse}. Venus de
l'eau~: une inondation. Venus du sol~: un \omterm{def:g5:natural-hazards:landslide}{glissement de terrain}, une
\omterm{def:g4:volcanoes-earthquakes-first-look:eruption}{éruption} volcanique, un \omterm{def:g4:volcanoes-earthquakes-first-look:earthquake}{séisme}.
\end{solution}

\begin{solution}{exo:g5:natural-hazards:3}
Une \omterm{def:g5:natural-hazards:drought}{sécheresse} est une longue période, de plusieurs mois ou de plusieurs
années, avec beaucoup moins de \omterm{def:g4:clouds-rain-snow:precipitation}{pluie} que d'habitude. Elle peut laisser les
agriculteurs sans récolte, parce que les cultures se fanent et meurent, et
elle peut laisser les villes à court d'eau potable, parce que les \omterm{def:g1:landscapes-around-us:water}{rivières}
et les \omterm{def:g1:landscapes-around-us:water}{lacs} rétrécissent.
\end{solution}

\begin{solution}{exo:g5:natural-hazards:4}
Par exemple~: rentrer ou attacher tout ce que le vent pourrait emporter
(chaises de jardin, jouets, poubelles)~; vérifier les lampes de poche et
garder prêtes des piles de rechange et de l'eau~; rester à l'intérieur,
loin des fenêtres, pendant la \omterm{def:g5:natural-hazards:storm}{tempête}~; continuer d'écouter les nouvelles
à la radio.
\end{solution}

\begin{solution}{exo:g5:natural-hazards:5}
(a) Un ouragan. (b) Un typhon. (c) Un cyclone tropical.
\end{solution}

\begin{solution}{exo:g5:natural-hazards:6}
Vert, jaune, orange, rouge. Au niveau rouge, le danger est grand~: il faut
agir tout de suite et suivre les consignes des personnes responsables.
\end{solution}

\begin{solution}{exo:g5:natural-hazards:7}
Une partie du versant de droite s'est détachée et a glissé vers le bas, en
laissant une cicatrice nue. Sa \omterm{def:g3:stones-pebbles-sand:sand}{boue} s'est étalée sur toute la largeur du
fond de la \omterm{def:g1:landscapes-around-us:landform}{vallée} et a recouvert la route. Elle a aussi barré la \omterm{def:g1:landscapes-around-us:water}{rivière}
comme un barrage~: l'eau de la \omterm{def:g1:landscapes-around-us:water}{rivière}, qui ne peut plus s'écouler, s'est
accumulée derrière et a inondé les bois en amont.
\end{solution}

\begin{solution}{exo:g5:natural-hazards:8}
Lena a raison. Nous ne pouvons pas arrêter la \omterm{def:g4:clouds-rain-snow:precipitation}{pluie} ni la \omterm{def:g1:landscapes-around-us:water}{rivière}, mais
nous pouvons mettre les gens à l'abri du danger. Par exemple~: construire
les maisons sur les hauteurs plutôt que dans la \omterm{def:g5:water-shapes-the-land:flood}{plaine d'inondation}~;
prévenir les gens à temps par des alertes aux inondations, pour qu'ils
montent sur les hauteurs~; apprendre à chacun à ne jamais traverser l'eau
d'une inondation, ni à pied ni en voiture.
\end{solution}

\begin{solution}{exo:g5:natural-hazards:9}
(a) $4 \times 60 = 240$ millimètres. (b) Le quart de 60, c'est 15~; en
quatre mois, $4 \times 15 = 60$ millimètres. (c) Il a manqué
$240 - 60 = 180$ millimètres~: les trois quarts de la \omterm{def:g4:clouds-rain-snow:precipitation}{pluie} habituelle.
\end{solution}

\begin{solution}{exo:g5:natural-hazards:10}
(a) $1\,200 \div 4 = 300$ logements. (b) Le tiers de 300, c'est 100~; il
reste donc $300 - 100 = 200$ logements dans la \omterm{def:g5:water-shapes-the-land:flood}{plaine d'inondation}.
\end{solution}

\begin{solution}{exo:g5:natural-hazards:11}
Raisons possibles~: le premier village avait été prévenu à temps et les gens
s'étaient préparés (volets fermés, objets rentrés), le second non~; les
maisons du premier sont solidement construites, celles du second sont
fragiles~; le second village se trouve sur une \omterm{def:g1:landscapes-around-us:water}{côte} basse ou dans une
\omterm{def:g5:water-shapes-the-land:flood}{plaine d'inondation}, où la \omterm{def:g5:natural-hazards:storm}{tempête} a aussi apporté une inondation, alors
que le premier est bâti sur une hauteur~; dans le second village, des
arbres proches des maisons sont tombés dessus. Le même aléa ne cause pas
les mêmes dégâts partout~: cela dépend de l'endroit où les gens vivent, de
leur façon de vivre et de leur degré de préparation.
\end{solution}
